% Overleaf-Datei (Name zu bestaetigen), Spiegel in docs/overleaf/.
% Nur zusammen mit der Overleaf-Datei aendern.
% Inhalt: Approval-Seite und Preface zum Umfang des Projekts.

\section*{Approval}
This thesis has been prepared over six months at the Department of
Energy Conversion and Storage, DTU Energy, at the Technical University
of Denmark, DTU, in partial fulfilment of the requirements for the
degree Master of Science in Engineering, MSc Eng.

It is assumed that the reader has a basic knowledge of quantum
chemistry and machine learning.

Artificial intelligence tools were used during the thesis work as
support tools. Claude (Anthropic), in the chat interface and through
Claude Code, was used for help with code development and debugging,
for the analysis and
plotting scripts, for checking the reported numbers against the
result files, for literature search, and for language correction and
the revision of the report text, which the author drafted and edited
throughout. The tools were used as assistance during the working
process, not as a replacement for independent analysis or scientific
judgement. All AI-generated suggestions were reviewed, edited and
checked before being included in the thesis. The calculations, the
modelling choices, the results, the discussion and the conclusions
are the responsibility of the author.

\vfill

\begin{center}
\namesigdate{\thesisauthor~-~\studentnumber}
\end{center}

\vfill

\newpage
\section*{Preface regarding the scope of this project}
\addcontentsline{toc}{section}{Preface}
This thesis was originally set up to answer only one of the two
research questions that will be explained in the introduction,
research question 1: whether a machine-learning potential trained on
Transition1x can be lifted to a higher level of theory by relabelling
a small part of its training data. Chapter~\ref{ch:delta} answers it.

The second question arose during that work. The OMol25 models were
used along the way as a point of comparison, and on some reactions
they did not find the restricted reference transition states but
different structures. The reason turned out to be that they are
trained on unrestricted labels, while the references we worked with
were optimised with restricted spin symmetry, which does not allow
broken-symmetry solutions. Where a transition state has one, the models
and our references describe different surfaces. Whether the OMol25
models can find the transition states of the unrestricted surface
there, although their training paths were never relaxed on it, became
research question 2.

Research question 2 was not part of the original scope. It was
followed for two reasons, and the change of direction was agreed with
the supervisors at every step. Research question 1 had meanwhile also
been addressed by others in related settings: relabelling with a small subset works.
Research question 2 had not been tested systematically, because the
published benchmarks score the models against restricted references.
We therefore wanted to address that gap. Chapter~\ref{ch:mr} is the
result.

This is only the short account; the introduction explains both
questions, and the opening of Chapter~\ref{ch:mr} explains how the
second follows from the first.
